% ===========================================================================
%  Resumen ejecutivo
%  Cuentas Nacionales de Transferencia para Mexico, 2016-2024
%
%  Compilar:  pdflatex resumen_ejecutivo   (sin bibliografia)
%  Las figuras las genera paper.ipynb; se copian a figures/ desde paper/figures.
% ===========================================================================
\documentclass[11pt,a4paper]{article}

\usepackage[utf8]{inputenc}
\usepackage[T1]{fontenc}
\usepackage[spanish,es-nodecimaldot]{babel}
\usepackage{lmodern}

\usepackage{graphicx}
\usepackage{xcolor}
\usepackage{caption}
\usepackage{float}
\usepackage{enumitem}
\usepackage[top=1.9cm,bottom=1.7cm,left=2.0cm,right=2.0cm]{geometry}
\usepackage[hidelinks]{hyperref}

\graphicspath{{figures/}}
\captionsetup{font=small,labelfont=bf}
\setlength{\parskip}{0.38em}
\setlength{\parindent}{0pt}
\pagestyle{empty}

\definecolor{acento}{HTML}{2C6FBB}

\newcommand{\titulillo}[1]{%
  \vspace{0.30em}\textcolor{acento}{\textbf{\large #1}}\par\vspace{0.08em}}

\begin{document}

\begin{center}
{\LARGE\bfseries Cuentas Nacionales de Transferencia para México}\\[0.35em]
{\large Quién consume, quién aporta y quién financia la diferencia, 2016--2024}\\[0.7em]
{\normalsize Juan Pablo López Reynosa \quad · \quad Héctor Villarreal}
\end{center}

\vspace{0.4em}
\hrule
\vspace{0.9em}

% ---------------------------------------------------------------------------
\titulillo{El ciclo de vida económico no se cierra igual para todos}

El consumo de una persona sólo coincide con su ingreso laboral en un tramo acotado de la edad
adulta. Esa diferencia ---el déficit del ciclo de vida--- se financia por dos vías: la privada,
con transferencias entre hogares y reasignación de activos del propio hogar, y la pública, con el
saldo entre lo que cada edad recibe del Estado y lo que le aporta. Cómo se reparte esa carga entre
ambas vías importa especialmente en una transición demográfica que envejecerá la composición
etaria del país. Este trabajo mide esos flujos año por año y edad por edad, con la Encuesta
Nacional de Ingresos y Gastos de los Hogares conciliada contra el Sistema de Cuentas Nacionales.

El resultado central es que ese ciclo se cierra de forma distinta para hombres y mujeres. El
perfil masculino completa el ciclo: entre los 33 y los 62 años el ingreso laboral de los hombres
supera a su consumo. El perfil femenino \textbf{no alcanza superávit en ninguna edad, en ninguno
de los cinco levantamientos analizados}. Las mujeres concentran 76.3\,\% del déficit agregado del
país.

El promedio nacional aparece deficitario en todas las edades, y conviene leerlo con precisión.
Agrega perfiles distintos entre sí, y su consumo incluye ---por convención del marco--- el
consumo de gobierno además del gasto de los hogares. Frente a ese consumo total, la participación
del ingreso laboral en el producto mexicano, entre 39\,\% y 44\,\% según el año, deja una
diferencia que existe en las cuentas nacionales antes del reparto por edad. El déficit
del promedio es un rasgo de la composición macroeconómica y de la agregación entre grupos, no un
juicio sobre una edad en particular.

\begin{figure}[H]
\centering
\includegraphics[width=0.88\linewidth]{fig_lcd_sexo.pdf}
\caption{Déficit del ciclo de vida por sexo. A la izquierda, el perfil por edad en 2024: el área
sombreada marca las edades en que el ingreso laboral masculino supera al consumo. A la derecha, el
saldo de la población de 25 a 59 años; los valores negativos son superávit.}
\end{figure}

% ---------------------------------------------------------------------------
\titulillo{El déficit no crece: se envejece}

El marco contable garantiza que exista déficit, de modo que la pregunta relevante no es si lo hay
sino cómo evoluciona y dónde se concentra. En términos reales el déficit agregado se mantuvo
prácticamente constante entre 2016 y 2024, y como proporción del consumo incluso descendió de
47.4\,\% a 42.2\,\%. Lo que sí cambió es su distribución por edad: el déficit por persona de 60
años o más creció 71.7\,\% en términos reales, mientras el peso de ese grupo en la población pasó
de 11.3\,\% a 15.1\,\%. El déficit es positivo y sostenido, y se está recorriendo hacia la vejez.

% ---------------------------------------------------------------------------
\titulillo{Quién financia a quién}

Los dos extremos del ciclo se financian por vías opuestas. El déficit de los menores de 15 años se
cubre en 88.6\,\% con recursos privados ---reasignación de activos del hogar y transferencias
entre hogares--- mientras el de los mayores de 60 se cubre en 54.5\,\% con transferencias públicas
netas. Un cambio demográfico que aumente el peso de la vejez se traduce de forma directa en
presión presupuestaria; uno que aumente el peso de la niñez se absorbe dentro de los hogares.

El saldo intergeneracional se movió de forma desigual. Entre 2016 y 2024 la transferencia pública
neta por adulto mayor subió 3.8\,\% real y el pago neto por persona en edad activa 39.1\,\%. Ese
aumento no viene de una mayor carga tributaria sino de lo que ese grupo dejó de recibir: su
participación en las transferencias públicas se redujo a la mitad, desplazada hacia las edades
mayores.

La desigualdad por sexo aparece también del lado del financiamiento. El déficit de las mujeres de
60 años o más duplica al de los hombres, y las transferencias públicas netas que reciben son
24.0\,\% menores: el sistema público no compensa la brecha, se suma a ella.

% ---------------------------------------------------------------------------
\titulillo{Qué aporta el procedimiento}

El objetivo del trabajo fue construir un procedimiento reproducible y escalable de Cuentas
Nacionales de Transferencia para México, que conecte la información de hogares con las cuentas
nacionales y sirva de base para cuentas de inclusión.

\begin{figure}[H]
\centering
\includegraphics[width=0.98\linewidth]{fig_pipeline.pdf}
\end{figure}

El procedimiento integra la encuesta de hogares y el sistema de cuentas nacionales para
identificar el ingreso laboral, el consumo, las transferencias que entran y salen de cada persona
---dentro del hogar, entre hogares y con el gobierno--- y la reasignación por activos. El
tratamiento del gobierno es deliberado: aparece como \textbf{mediador entre generaciones y no como
agente}, porque recauda de unas edades y transfiere a otras sin consumir por cuenta propia dentro
de la identidad.

Cada variable se construye con una decisión metodológica explícita, documentada y sustituible. Esa
es la propiedad que hace escalable al procedimiento: mejorar la identificación de un bloque no
exige rehacer los demás. Un ejemplo ilustra el criterio. La reasignación por activos se ancla a un
agregado neto de ahorro, lo que permite cuadrar el déficit por sus dos lados sin repartir el
ahorro por edad. Es una decisión válida para el propósito actual y no pretende ser un instrumento
para medir desigualdad del ingreso de capital, que requeriría otro diseño.

% ---------------------------------------------------------------------------
\titulillo{Siguiente paso}

El producto es una base con observaciones a nivel de agente representativo, construida desde la
encuesta y conciliada con las cuentas nacionales. Esa estructura habilita el paso siguiente:
proyectar los perfiles bajo distintos escenarios demográficos y económicos, y evaluar el efecto de
cambios de política sobre cada grupo de edad y sexo.

\vspace{0.6em}
\hrule
\vspace{0.4em}
{\footnotesize Cifras en pesos constantes de 2024. Fuente: elaboración propia con ENIGH (INEGI) y
SCNM. El documento metodológico detalla imputaciones, agregados de control y limitaciones.}

\end{document}
